\section*{September 4, 2026: Reassessing the school and housing comparison}
\textit{Agent-drafted research entry; recommendations remain open for discussion.}

Homes sold near the selected closed and still-open schools differ
substantially in price and building size before closure. Jacob requested
this review of the 30-closed/49-open-school comparison. The selected sample
contains 1,374 treated and 3,157 control
sales in 2008--2012, assigned to 29 treated and 48 control school sites.
Treated sales average \$218,206, compared with \$161,131 for controls; median
prices are \$127,432 and \$69,184. Mean building sizes are 2,342 and 2,039
square feet. Prices are in 2022 dollars.

A descriptive calculation first averaging sales within each site over
2008--2012, then weighting sites equally, gives means of \$205,132 for treated
sites and \$185,908 for controls. This smaller gap shows that the mix of sites
contributing sales matters. It is not a new DiD estimate. Neither weighting
scheme by itself establishes comparability. The current regressions give
each transaction equal weight; they do not estimate a simple equally weighted
average of school effects. Changes in which properties sell can also change
the measured price outcome. Housing controls address observed differences
among sales but cannot recover unobserved prices for properties that do not sell.

The source school ledger defines a narrower treatment than all 2013 closures.
It flags 47 closed schools, but excludes 17 whose buildings remained in active
school use, leaving 30 treated schools at 29 physical sites. The 49 control
schools are also a selected subset of February candidates that remained open.
These definitions are inherited from the supplied ledger. School closure and
loss of school use at a building are different treatments. A treatment based
on the adopted closure plan differs from selecting sites by subsequently
realized vacancy or reuse; the ledger notes alone do not establish the
information date of every classification.

The quarter-mile sample additionally excludes proximity to welcoming schools
and other candidate sites. Among pre-closure sales within a quarter mile of one treatment group
and outside that radius of the other, this removes 461 of 1,835 treated sales,
compared with 101 of 3,258
control sales. These exclusions address competing exposure but also change
the geographic comparison. These estimates concern homes near the school buildings. Distance alone
does not establish which school children at those addresses attended.

Codex recommends retaining the candidate list as the initial comparison pool
while examining overlap at the school-neighborhood level using pre-treatment
housing, geography, school conditions, and neighborhood recovery measures.
Matching or weighting should use an explicit baseline rule and report sites
without credible counterparts. A nearby outer-ring comparison could provide
a useful second design, but estimates a relative local effect and requires
attention to spillovers. No treatment definition, sample, weights, or
regression was changed in this review.

{\footnotesize Reproduction: the companion calculation prints the two price
weighting comparisons from the TWFE audit's site/year support, using
Make-managed input links. Other descriptive values come from the exposure
audit's treatment/control summary. School definitions and notes are in the
raw-school task's supplied closure ledger, reproduced in the cleaned school
file. These are the September 4 non-condo, annual-price-per-square-foot-trimmed
outputs at revision c200c31 plus working changes. The source notes were
reviewed; underlying closure-plan documents were not freshly verified.}
